
\begin{exercise}[Liminf and Limsup Properties \Coffeecup]
    \label{ex: liminf and limsup properties}
    Prove Lemma \ref{lem: liminf and limsup properties}. 

    \begin{solution}
        \begin{enumerate}[label=(\roman*)]
            \item To show that  \(\liminf_n A_n \subseteq \limsup_n A_n\) let
            \[
                \omega \in \liminf_n A_n = \bigcup_{N\in \nat} \bigcap_{n\ge N} A_n.
            \]
            then there exists an \(N\) such that for all \(n \ge N\) we have
            \(\omega \in A_n\). In particular, there are infinitely many \(n\)
            such that \(\omega \in A_n\) and consequently for all \(N\in \nat\) there
            exists \(n\ge N\) such that \(\omega \in A_n\), which implies
            \[
                \omega \in \bigcap_{N\in \nat} \bigcup_{n\ge N} A_n = \limsup_n A_n.
            \]
            
            \item\label{it: de morgans on liminf/limsup} We have using De Morgan's laws that
            \[
            \begin{aligned}
                (\limsup_n A_n)^\complement
                &= \Bigl(\bigcap_{N\in \nat} \bigcup_{n\ge N} A_n\Bigr)^\complement
                = \bigcup_{N\in \nat} \Bigl(\bigcup_{n\ge N} A_n\Bigr)^\complement
                = \bigcup_{N\in\nat} \bigcap_{n\ge N} A_n^\complement
                \\
                &= \liminf_n A_n^\complement.
            \end{aligned}
            \]

            \item\label{it: function limsup to set limsup} To prove \(\limsup_n \ind{A_n} = \ind*{\limsup_n A_n}\) 
            observe that 
            \[
                \limsup_n \ind{A_n}(\omega)
                = \lim_{N\to \infty} \sup_{n\ge N} \ind{A_n}(\omega)
            \]
            is in the set \(\set{0,1}\) for all \(\omega\) since
            \(\ind{A_n}(\omega) \in \set{0,1}\) for all \(n\) and
            the supremum is thereby either \(0\) or \(1\). We therefore only need to show
            that \(\limsup_n \ind{A_n}(\omega) = 1\) if and only if
            \begin{equation}
                \label{eq: omega in limsup_n An} 
                \omega \in \limsup_n A_n = \bigcap_{N\in \nat} \bigcup_{n\ge N} A_n.
            \end{equation}
            Assume \eqref{eq: omega in limsup_n An} holds, then for all \(N\in
            \nat\) there exists an \(n \ge N\) such that \(\omega \in A_n\). In particular,
            since for any \(N\in \nat\) we can find \(n\ge N\) such that \(\ind{A_n}(\omega) = 1\)
            and therefore \(\sup_{n\ge N} \ind{A_n}(\omega) = 1\). Conversely, if
            \(\limsup_n \ind{A_n}(\omega) = 1\), then for all \(N\in \nat\) we have
            \(\sup_{n\ge N} \ind{A_n}(\omega) =1\) and consequently there must exist
            \(n\ge N\) such that \(\ind{A_n}(\omega) = 1\) and therefore
            \(\omega \in A_n\). This implies \eqref{eq: omega in limsup_n An}.

            \item Observe that we can translate this to the previous results as follows:
            \[\begin{aligned}[b]
                \liminf_n \ind{A_n}
                &= \liminf_n (1-\ind{A_n^\complement})
                \\
                &= 1 - \limsup_n \ind{A_n^\complement}
                \\
                \overset{\text{\ref{it: function limsup to set limsup}}}&= 1 - \ind*{\limsup_n A_n^\complement}
                \overset{\text{\ref{it: de morgans on liminf/limsup}}}= \ind*{\liminf_n A_n}.
            \end{aligned}
            \qedhere
            \]
        \end{enumerate}        
    \end{solution}
\end{exercise}
